\documentclass[11pt,a4paper]{article}
\usepackage[utf8]{inputenc}
\usepackage[T1]{fontenc}
\usepackage[margin=1in]{geometry}
\usepackage{graphicx}
\usepackage{hyperref}
\usepackage{xcolor}
\usepackage{float}
\usepackage{listings}
\usepackage{upquote}
\usepackage{amsmath}
\usepackage{enumitem}
\usepackage{tcolorbox}
\usepackage{booktabs}
\tcbuselibrary{skins}

% Code listing style
\lstset{
    basicstyle=\ttfamily\small,
    breaklines=true,
    breakatwhitespace=false,
    frame=single,
    numbers=left,
    numberstyle=\tiny,
    backgroundcolor=\color{gray!10},
    keywordstyle=\color{blue},
    commentstyle=\color{green!50!black},
    stringstyle=\color{red},
    showstringspaces=false
}

% Callout boxes
\newtcolorbox{notebox}{colback=blue!5,colframe=blue!40!black,title=Note,fonttitle=\bfseries}
\newtcolorbox{bonusbox}{colback=green!5,colframe=green!40!black,title=Bonus Opportunity,fonttitle=\bfseries}

\title{\textbf{Final Project --- Stretch Robot Proof of Concept}}
\author{
    Robotics and Cyber-Physical Systems \\
    School of Sciences and Engineering \\
    Tecnol\'ogico de Monterrey
}
\date{}

\begin{document}

\maketitle

\section*{Overview}

For the final project, your team will identify a practical use case for the Hello Robot Stretch and implement a proof-of-concept demonstration that shows the idea is feasible given the robot's hardware.

The goal is not a finished, production-ready system. It is a convincing demonstration: something that works well enough to make a clear case that the application could be built out further. Think of it as pitching an idea to a potential collaborator or investor---you want to leave them believing the concept is real and achievable.

The \texttt{stretch\_mujoco\_digital\_twin} simulator is your primary development platform. Scripts written against \texttt{stretch\_toolkit} are designed to port to the physical Stretch robot with minimal changes.

\section*{Choosing a Use Case}

There are no restrictions on the application domain. Any use case that makes genuine use of the Stretch robot's capabilities is acceptable. Some directions to consider:

\begin{itemize}
    \item Assistive robotics (helping people with limited mobility perform everyday tasks)
    \item Object manipulation and sorting
    \item Inspection or monitoring tasks
    \item Human--robot interaction scenarios
    \item Navigation and fetch tasks
\end{itemize}

Your use case should be specific enough that you can scope a meaningful proof of concept within the project timeline. Broad ideas (``a general-purpose home robot'') are harder to demonstrate convincingly than focused ones (``a robot that retrieves a specific object from a shelf on request'').

\section*{Deliverables}

Every team must submit the following:

\begin{enumerate}
    \item \textbf{Source code} --- A working script (or set of scripts) that runs in the simulator. Code should be reasonably clean and include brief comments explaining non-obvious sections. Submit via the course platform or as a link to a repository.

    \item \textbf{Written report} (2--3 pages) --- A concise document covering:
    \begin{itemize}
        \item The use case and why the Stretch robot is a suitable platform for it
        \item A description of how the system works
        \item Known limitations and what would be needed to take it further
    \end{itemize}

    \item \textbf{Video} --- A recording that shows your demo running. This can be a raw screen capture or a narrated/edited video---whatever best conveys what the system does. The video should be self-explanatory to someone who has not seen your presentation.
\end{enumerate}

\textbf{File naming:} \texttt{Team\_[TeamName]\_FinalProject.zip} containing all of the above.

\section*{Presentation}

Each team will present to the class. The presentation should include:

\begin{itemize}
    \item \textbf{Slides} --- Cover the use case, your approach, and the key design decisions.
    \item \textbf{Demo} --- Show your system running live in the simulator. If a live demo is not practical on the day, a pre-recorded video of the simulation is acceptable as a fallback.
    \item \textbf{Q\&A} --- Be prepared to answer questions about your implementation and design choices.
\end{itemize}

Presentations should be roughly 10--15 minutes including Q\&A.

\begin{bonusbox}
Teams that successfully run their proof-of-concept on the \textbf{physical Stretch robot} are exempt from the class presentation. You still need to submit all three written deliverables (code, report, video).

To arrange a session, use the scheduling tool provided by the instructor to book a time slot. Not every team member needs to attend---one or two representatives are sufficient. Multiple teams are welcome to coordinate and share a single session, taking turns with the robot.
\end{bonusbox}

\section*{Grading}

The project is graded across four equally-weighted areas:

\begin{center}
\begin{tabular}{llp{8.5cm}}
\toprule
\textbf{Area} & \textbf{Weight} & \textbf{What we are looking for} \\
\midrule
\textbf{Idea \& Feasibility} & 25\% &
    Does the use case make sense for this robot? Is the scope appropriate---ambitious enough to be interesting, bounded enough to demonstrate? \\[6pt]
\textbf{Implementation} & 25\% &
    Does the demo work? Does the code make reasonable use of the robot's capabilities? Is there some autonomous behavior, or is it purely teleoperated? \\[6pt]
\textbf{Presentation} & 25\% &
    Is the pitch clear and convincing? Does the demo (live or recorded) support the argument? Can the team answer questions about their system? \\[6pt]
\textbf{Video} & 25\% &
    Does the video clearly show what the system does? Would someone watching it understand the idea without additional explanation? \\
\bottomrule
\end{tabular}
\end{center}

\begin{notebox}
These criteria are intentionally broad. There is no single correct way to build a proof of concept, and grading will take into account the difficulty and ambition of what you attempted relative to what you achieved.
\end{notebox}

\end{document}
